%%%%%%%%%%%%%%%%%%%%%%%%%%%%%%%%%%%%%%%%%%%%%%%%%%%%%%%%%%%%%%%%%%%%%%%%%%%%%%%%
% CHAPTER 9
%%%%%%%%%%%%%%%%%%%%%%%%%%%%%%%%%%%%%%%%%%%%%%%%%%%%%%%%%%%%%%%%%%%%%%%%%%%%%%%%

\chapter{Discrete Fourier Transform}

\begin{tcolorbox}[colback=blue!5!white,colframe=blue!70!black,title=Learning Objectives]

After studying this chapter, the reader will be able to

\begin{itemize}
\item Understand the need for the DFT.
\item Derive the DFT and inverse DFT.
\item Understand the relationship between DTFT and DFT.
\item Compute DFTs of finite sequences.
\item Understand twiddle factors and their properties.
\item Interpret DFT bins and frequency resolution.
\end{itemize}

\end{tcolorbox}

%%%%%%%%%%%%%%%%%%%%%%%%%%%%%%%%%%%%%%%%%%%%%%%%%%%%%%%%%%%%%%%%%%%%%%%%%%%%%%%%

\section{Introduction}

The DTFT of a sequence is continuous in frequency and contains infinitely many
frequency samples. A digital computer cannot store or process an infinite number
of frequency values.

For a finite sequence of length \(N\), we compute the DTFT only at \(N\)
equally spaced frequencies. These samples form the
\textbf{Discrete Fourier Transform (DFT)}.

%%%%%%%%%%%%%%%%%%%%%%%%%%%%%%%%%%%%%%%%%%%%%%%%%%%%%%%%%%%%%%%%%%%%%%%%%%%%%%%%

\section{From DTFT to DFT}

For a finite sequence \(x[n]\) of length \(N\),

\[
X(e^{j\omega})=
\sum_{n=0}^{N-1}x[n]e^{-j\omega n}.
\]

Evaluate it at

\[
\omega_k=\frac{2\pi k}{N},
\qquad k=0,1,\ldots,N-1.
\]

Then

\[
X[k]=X(e^{j\omega_k}).
\]

%%%%%%%%%%%%%%%%%%%%%%%%%%%%%%%%%%%%%%%%%%%%%%%%%%%%%%%%%%%%%%%%%%%%%%%%%%%%%%%%

\section{DFT Definition}

\[
\boxed{
X[k]=
\sum_{n=0}^{N-1}
x[n]e^{-j\frac{2\pi}{N}kn}
}
\]

where

\[
k=0,1,\ldots,N-1.
\]

Here,

\begin{itemize}
\item \(x[n]\): time-domain sequence,
\item \(X[k]\): DFT coefficients,
\item \(N\): number of samples.
\end{itemize}

%%%%%%%%%%%%%%%%%%%%%%%%%%%%%%%%%%%%%%%%%%%%%%%%%%%%%%%%%%%%%%%%%%%%%%%%%%%%%%%%

\section{Inverse DFT}

The original sequence is recovered by

\[
\boxed{
x[n]=
\frac{1}{N}
\sum_{k=0}^{N-1}
X[k]e^{j\frac{2\pi}{N}kn}
}
\]

Thus, the DFT and IDFT form a transform pair.

%%%%%%%%%%%%%%%%%%%%%%%%%%%%%%%%%%%%%%%%%%%%%%%%%%%%%%%%%%%%%%%%%%%%%%%%%%%%%%%%

\section{Twiddle Factor}

Define

\[
\boxed{
W_N=e^{-j\frac{2\pi}{N}}
}
\]

Then the DFT becomes

\[
\boxed{
X[k]=
\sum_{n=0}^{N-1}
x[n]W_N^{kn}
}
\]

and the IDFT becomes

\[
\boxed{
x[n]=
\frac{1}{N}
\sum_{k=0}^{N-1}
X[k]W_N^{-kn}.
}
\]

%%%%%%%%%%%%%%%%%%%%%%%%%%%%%%%%%%%%%%%%%%%%%%%%%%%%%%%%%%%%%%%%%%%%%%%%%%%%%%%%

\section{Properties of the Twiddle Factor}

%%%%%%%%%%%%%%%%%%%%%%%%%%%%%%%%%%%%%%%%%%%%%%%%%%%%%%%%%%%%%%%%%%%%%%%%%%%%%%%%

\subsection{Periodicity}

\[
\boxed{
W_N^{k+N}=W_N^k
}
\]

%%%%%%%%%%%%%%%%%%%%%%%%%%%%%%%%%%%%%%%%%%%%%%%%%%%%%%%%%%%%%%%%%%%%%%%%%%%%%%%%

\subsection{Symmetry}

\[
\boxed{
W_N^{k+N/2}=-W_N^k
}
\]

for even \(N\).

%%%%%%%%%%%%%%%%%%%%%%%%%%%%%%%%%%%%%%%%%%%%%%%%%%%%%%%%%%%%%%%%%%%%%%%%%%%%%%%%

\subsection{Complex Conjugate}

\[
\boxed{
(W_N^k)^*=W_N^{-k}
}
\]

These properties are heavily used in FFT algorithms.

%%%%%%%%%%%%%%%%%%%%%%%%%%%%%%%%%%%%%%%%%%%%%%%%%%%%%%%%%%%%%%%%%%%%%%%%%%%%%%%%

\section{Orthogonality}

The complex exponentials are orthogonal:

\[
\boxed{
\sum_{n=0}^{N-1}W_N^{(k-m)n}=
\begin{cases}
N, & k=m\\
0, & k\neq m
\end{cases}
}
\]

Orthogonality guarantees perfect reconstruction.

%%%%%%%%%%%%%%%%%%%%%%%%%%%%%%%%%%%%%%%%%%%%%%%%%%%%%%%%%%%%%%%%%%%%%%%%%%%%%%%%

\section{Matrix Representation}

Define the DFT matrix

\[
[F_N]_{k,n}=W_N^{kn}.
\]

Then

\[
\boxed{
\mathbf{X}=F_N\mathbf{x}
}
\]

and

\[
\boxed{
\mathbf{x}=\frac{1}{N}F_N^H\mathbf{X}
}
\]

where \(F_N^H\) is the conjugate transpose of \(F_N\).

%%%%%%%%%%%%%%%%%%%%%%%%%%%%%%%%%%%%%%%%%%%%%%%%%%%%%%%%%%%%%%%%%%%%%%%%%%%%%%%%

\section{Frequency Bins}

The \(k\)-th DFT coefficient corresponds to frequency

\[
\boxed{
f_k=\frac{k}{N}f_s
}
\]

where \(f_s\) is the sampling frequency.

%%%%%%%%%%%%%%%%%%%%%%%%%%%%%%%%%%%%%%%%%%%%%%%%%%%%%%%%%%%%%%%%%%%%%%%%%%%%%%%%

\subsection{Frequency Resolution}

The spacing between adjacent bins is

\[
\boxed{
\Delta f=\frac{f_s}{N}
}
\]

Larger \(N\) gives finer frequency resolution.

%%%%%%%%%%%%%%%%%%%%%%%%%%%%%%%%%%%%%%%%%%%%%%%%%%%%%%%%%%%%%%%%%%%%%%%%%%%%%%%%

\section{Example 1: 2-Point DFT}

Let

\[
x[0]=1,\qquad x[1]=1.
\]

For \(N=2\),

\[
W_2=e^{-j\pi}=-1.
\]

%%%%%%%%%%%%%%%%%%%%%%%%%%%%%%%%%%%%%%%%%%%%%%%%%%%%%%%%%%%%%%%%%%%%%%%%%%%%%%%%

\subsection{Compute \(X[0]\)}

\[
X[0]=1+1=2.
\]

%%%%%%%%%%%%%%%%%%%%%%%%%%%%%%%%%%%%%%%%%%%%%%%%%%%%%%%%%%%%%%%%%%%%%%%%%%%%%%%%

\subsection{Compute \(X[1]\)}

\[
X[1]=1+1(-1)=0.
\]

%%%%%%%%%%%%%%%%%%%%%%%%%%%%%%%%%%%%%%%%%%%%%%%%%%%%%%%%%%%%%%%%%%%%%%%%%%%%%%%%

\subsection{Result}

\[
\boxed{
X[k]=\{2,0\}
}
\]

The sequence contains only a DC component.

%%%%%%%%%%%%%%%%%%%%%%%%%%%%%%%%%%%%%%%%%%%%%%%%%%%%%%%%%%%%%%%%%%%%%%%%%%%%%%%%

\section{Example 2: 4-Point DFT}

Let

\[
x[n]=\{1,0,0,0\}.
\]

Using the DFT definition,

\[
X[k]=1
\]

for all \(k\).

Hence,

\[
\boxed{
X[k]=\{1,1,1,1\}
}
\]

This is the discrete-time counterpart of

\[
\delta[n]\leftrightarrow1.
\]

%%%%%%%%%%%%%%%%%%%%%%%%%%%%%%%%%%%%%%%%%%%%%%%%%%%%%%%%%%%%%%%%%%%%%%%%%%%%%%%%

\section{Example 3: Constant Sequence}

Let

\[
x[n]=\{1,1,1,1\}.
\]

%%%%%%%%%%%%%%%%%%%%%%%%%%%%%%%%%%%%%%%%%%%%%%%%%%%%%%%%%%%%%%%%%%%%%%%%%%%%%%%%

\subsection{Compute \(X[0]\)}

\[
X[0]=1+1+1+1=4.
\]

%%%%%%%%%%%%%%%%%%%%%%%%%%%%%%%%%%%%%%%%%%%%%%%%%%%%%%%%%%%%%%%%%%%%%%%%%%%%%%%%

\subsection{Compute \(X[1]\)}

\[
X[1]=1+W_4+W_4^2+W_4^3.
\]

Since

\[
W_4=e^{-j\pi/2}=-j,
\]

the terms cancel:

\[
X[1]=0.
\]

Similarly,

\[
X[2]=0,\qquad X[3]=0.
\]

%%%%%%%%%%%%%%%%%%%%%%%%%%%%%%%%%%%%%%%%%%%%%%%%%%%%%%%%%%%%%%%%%%%%%%%%%%%%%%%%

\subsection{Result}

\[
\boxed{
X[k]=\{4,0,0,0\}
}
\]

Only the DC bin is nonzero.

%%%%%%%%%%%%%%%%%%%%%%%%%%%%%%%%%%%%%%%%%%%%%%%%%%%%%%%%%%%%%%%%%%%%%%%%%%%%%%%%

\section{DFT of a Complex Exponential}

Let

\[
x[n]=e^{j\frac{2\pi}{N}mn}.
\]

Then

\[
X[k]=
\sum_{n=0}^{N-1}
e^{j\frac{2\pi}{N}(m-k)n}.
\]

Using orthogonality,

\[
\boxed{
X[k]=
\begin{cases}
N, & k=m\\
0, & k\neq m
\end{cases}
}
\]

Thus, a complex exponential occupies exactly one DFT bin.

%%%%%%%%%%%%%%%%%%%%%%%%%%%%%%%%%%%%%%%%%%%%%%%%%%%%%%%%%%%%%%%%%%%%%%%%%%%%%%%%

\section{DFT Spectrum Interpretation}

For a real sequence:

\begin{itemize}
\item \(X[0]\): DC component,
\item \(X[1]\) to \(X[N/2-1]\): positive frequencies,
\item \(X[N/2]\): Nyquist component (for even \(N\)),
\item remaining bins: negative frequencies.
\end{itemize}

%%%%%%%%%%%%%%%%%%%%%%%%%%%%%%%%%%%%%%%%%%%%%%%%%%%%%%%%%%%%%%%%%%%%%%%%%%%%%%%%

\section{Conjugate Symmetry}

If \(x[n]\) is real,

\[
\boxed{
X[N-k]=X^*[k]
}
\]

Hence,

\[
|X[N-k]|=|X[k]|
\]

and

\[
\angle X[N-k]=-\angle X[k].
\]

Only half the spectrum is independent.

%%%%%%%%%%%%%%%%%%%%%%%%%%%%%%%%%%%%%%%%%%%%%%%%%%%%%%%%%%%%%%%%%%%%%%%%%%%%%%%%

\section{Example 4: Real Sequence}

Let

\[
x[n]=\{1,2,1,0\}.
\]

The DFT satisfies

\[
X[3]=X^*[1].
\]

So after computing \(X[1]\), \(X[3]\) is obtained automatically.

%%%%%%%%%%%%%%%%%%%%%%%%%%%%%%%%%%%%%%%%%%%%%%%%%%%%%%%%%%%%%%%%%%%%%%%%%%%%%%%%

\section{Circular Periodicity}

The DFT assumes that the sequence is periodic with period \(N\):

\[
\boxed{
x[n+N]=x[n]
}
\]

This assumption leads to \textbf{circular convolution}, which will be studied in
the next part.

%%%%%%%%%%%%%%%%%%%%%%%%%%%%%%%%%%%%%%%%%%%%%%%%%%%%%%%%%%%%%%%%%%%%%%%%%%%%%%%%

\section{Relationship Between DTFT and DFT}

\[
\boxed{
X[k]=X(e^{j\omega})\Big|_{\omega=2\pi k/N}
}
\]

Therefore, the DFT is simply a set of samples of the DTFT.

Increasing \(N\) gives more samples of the DTFT and a better approximation of the
continuous spectrum.

%%%%%%%%%%%%%%%%%%%%%%%%%%%%%%%%%%%%%%%%%%%%%%%%%%%%%%%%%%%%%%%%%%%%%%%%%%%%%%%%

\section{Zero Padding}

If a sequence has length \(L < N\), append zeros:

\[
x[n]=\{x_0,x_1,\ldots,x_{L-1},0,\ldots,0\}.
\]

%%%%%%%%%%%%%%%%%%%%%%%%%%%%%%%%%%%%%%%%%%%%%%%%%%%%%%%%%%%%%%%%%%%%%%%%%%%%%%%%

\subsection{Effect}

\begin{itemize}
\item No new frequency content is added.
\item The DTFT is sampled more densely.
\item Spectral plots appear smoother.
\end{itemize}

%%%%%%%%%%%%%%%%%%%%%%%%%%%%%%%%%%%%%%%%%%%%%%%%%%%%%%%%%%%%%%%%%%%%%%%%%%%%%%%%

\section{Quick Check}

For \(N=8\),

\[
\Delta f=\frac{f_s}{8}.
\]

If \(f_s=8\text{ kHz}\),

\[
\boxed{
\Delta f=1\text{ kHz}
}
\]

The DFT bins correspond to

\[
0,1,2,\ldots,7\text{ kHz}.
\]

%%%%%%%%%%%%%%%%%%%%%%%%%%%%%%%%%%%%%%%%%%%%%%%%%%%%%%%%%%%%%%%%%%%%%%%%%%%%%%%%

\begin{tcolorbox}[title=One-Minute Revision,colback=blue!5]

\[
X[k]=
\sum_{n=0}^{N-1}
x[n]e^{-j\frac{2\pi}{N}kn}
\]

\[
x[n]=
\frac{1}{N}
\sum_{k=0}^{N-1}
X[k]e^{j\frac{2\pi}{N}kn}
\]

\[
W_N=e^{-j\frac{2\pi}{N}}
\]

\[
f_k=\frac{k}{N}f_s
\]

\[
\Delta f=\frac{f_s}{N}
\]

Real sequence:

\[
X[N-k]=X^*[k]
\]

Complex exponential:

\[
e^{j2\pi mn/N}\rightarrow
X[m]=N
\]

DFT = sampled DTFT.

\end{tcolorbox}
%%%%%%%%%%%%%%%%%%%%%%%%%%%%%%%%%%%%%%%%%%%%%%%%%%%%%%%%%%%%%%%%%%%%%%%%%%%%%%%%
%%%%%%%%%%%%%%%%%%%%%%%%%%%%%%%%%%%%%%%%%%%%%%%%%%%%%%%%%%%%%%%%%%%%%%%%%%%%%%%%
\section{Linearity Property}

If

\[
x_1[n]\xleftrightarrow{\text{DFT}}X_1[k]
\]

and

\[
x_2[n]\xleftrightarrow{\text{DFT}}X_2[k],
\]

then

\[
\boxed{
a_1x_1[n]+a_2x_2[n]
\xleftrightarrow{\text{DFT}}
a_1X_1[k]+a_2X_2[k]
}
\]

%%%%%%%%%%%%%%%%%%%%%%%%%%%%%%%%%%%%%%%%%%%%%%%%%%%%%%%%%%%%%%%%%%%%%%%%%%%%%%%%

\section{Circular Time Shift}

A circular delay of \(n_0\) samples gives

\[
\boxed{
x[(n-n_0)_N]
\xleftrightarrow{\text{DFT}}
W_N^{kn_0}X[k]
}
\]

where \((\cdot)_N\) denotes modulo-\(N\) indexing.

%%%%%%%%%%%%%%%%%%%%%%%%%%%%%%%%%%%%%%%%%%%%%%%%%%%%%%%%%%%%%%%%%%%%%%%%%%%%%%%%

\section{Circular Frequency Shift}

Multiplying by a complex exponential shifts the DFT bins:

\[
\boxed{
x[n]W_N^{-mn}
\xleftrightarrow{\text{DFT}}
X[(k-m)_N]
}
\]

%%%%%%%%%%%%%%%%%%%%%%%%%%%%%%%%%%%%%%%%%%%%%%%%%%%%%%%%%%%%%%%%%%%%%%%%%%%%%%%%

\section{Circular Convolution}

%%%%%%%%%%%%%%%%%%%%%%%%%%%%%%%%%%%%%%%%%%%%%%%%%%%%%%%%%%%%%%%%%%%%%%%%%%%%%%%%

\subsection{Definition}

For two length-\(N\) sequences,

\[
\boxed{
y[n]=x[n]\circledast h[n]
}
\]

where

\[
\boxed{
y[n]=
\sum_{m=0}^{N-1}
x[m]h[(n-m)_N]
}
\]

The index is taken modulo \(N\).

%%%%%%%%%%%%%%%%%%%%%%%%%%%%%%%%%%%%%%%%%%%%%%%%%%%%%%%%%%%%%%%%%%%%%%%%%%%%%%%%

\section{Circular Convolution Property}

If

\[
x[n]\xleftrightarrow{\text{DFT}}X[k]
\]

and

\[
h[n]\xleftrightarrow{\text{DFT}}H[k],
\]

then

\[
\boxed{
x[n]\circledast h[n]
\xleftrightarrow{\text{DFT}}
X[k]H[k]
}
\]

This is the most important property of the DFT.

%%%%%%%%%%%%%%%%%%%%%%%%%%%%%%%%%%%%%%%%%%%%%%%%%%%%%%%%%%%%%%%%%%%%%%%%%%%%%%%%

\section{Example: Circular Convolution}

Let

\[
x[n]=\{1,2\},
\qquad
h[n]=\{1,1\}.
\]

Assume \(N=2\).

%%%%%%%%%%%%%%%%%%%%%%%%%%%%%%%%%%%%%%%%%%%%%%%%%%%%%%%%%%%%%%%%%%%%%%%%%%%%%%%%

\subsection{Compute \(y[0]\)}

\[
y[0]=x[0]h[0]+x[1]h[1]
\]

\[
=1(1)+2(1)=3.
\]

%%%%%%%%%%%%%%%%%%%%%%%%%%%%%%%%%%%%%%%%%%%%%%%%%%%%%%%%%%%%%%%%%%%%%%%%%%%%%%%%

\subsection{Compute \(y[1]\)}

\[
y[1]=x[0]h[1]+x[1]h[0]
\]

\[
=1(1)+2(1)=3.
\]

%%%%%%%%%%%%%%%%%%%%%%%%%%%%%%%%%%%%%%%%%%%%%%%%%%%%%%%%%%%%%%%%%%%%%%%%%%%%%%%%

\subsection{Result}

\[
\boxed{
y[n]=\{3,3\}
}
\]

%%%%%%%%%%%%%%%%%%%%%%%%%%%%%%%%%%%%%%%%%%%%%%%%%%%%%%%%%%%%%%%%%%%%%%%%%%%%%%%%

\section{Linear Convolution}

Ordinary convolution is

\[
\boxed{
y[n]=x[n]*h[n]
}
\]

with output length

\[
\boxed{
L=L_x+L_h-1
}
\]

Circular convolution and linear convolution are generally different.

%%%%%%%%%%%%%%%%%%%%%%%%%%%%%%%%%%%%%%%%%%%%%%%%%%%%%%%%%%%%%%%%%%%%%%%%%%%%%%%%

\section{Linear Convolution Using DFT}

To obtain linear convolution using DFT:

\begin{enumerate}
\item Choose
\[
N\ge L_x+L_h-1.
\]

\item Zero-pad both sequences to length \(N\).

\item Compute the DFTs.

\item Multiply:
\[
Y[k]=X[k]H[k].
\]

\item Compute the IDFT of \(Y[k]\).
\end{enumerate}

%%%%%%%%%%%%%%%%%%%%%%%%%%%%%%%%%%%%%%%%%%%%%%%%%%%%%%%%%%%%%%%%%%%%%%%%%%%%%%%%

\section{Example: Linear Convolution Using DFT}

Let

\[
x[n]=\{1,2\},
\qquad
h[n]=\{1,1\}.
\]

Output length:

\[
L=2+2-1=3.
\]

Choose \(N=3\).

%%%%%%%%%%%%%%%%%%%%%%%%%%%%%%%%%%%%%%%%%%%%%%%%%%%%%%%%%%%%%%%%%%%%%%%%%%%%%%%%

\subsection{Zero Padding}

\[
x[n]=\{1,2,0\},
\qquad
h[n]=\{1,1,0\}.
\]

%%%%%%%%%%%%%%%%%%%%%%%%%%%%%%%%%%%%%%%%%%%%%%%%%%%%%%%%%%%%%%%%%%%%%%%%%%%%%%%%

\subsection{Result}

The IDFT gives

\[
\boxed{
y[n]=\{1,3,2\}
}
\]

which matches the ordinary convolution result.

%%%%%%%%%%%%%%%%%%%%%%%%%%%%%%%%%%%%%%%%%%%%%%%%%%%%%%%%%%%%%%%%%%%%%%%%%%%%%%%%

\section{Circular vs Linear Convolution}

\begin{table}[H]
\centering
\caption{Comparison of convolution types}
\begin{tabular}{|p{5cm}|p{4cm}|p{4cm}|}
\hline
\textbf{Feature} & \textbf{Linear} & \textbf{Circular}\\
\hline
Output length & \(L_x+L_h-1\) & \(N\)\\
\hline
Indexing & Ordinary & Modulo \(N\)\\
\hline
Used in LTI systems & Yes & No (directly)\\
\hline
Used in DFT processing & Via zero padding & Yes\\
\hline
\end{tabular}
\end{table}

%%%%%%%%%%%%%%%%%%%%%%%%%%%%%%%%%%%%%%%%%%%%%%%%%%%%%%%%%%%%%%%%%%%%%%%%%%%%%%%%

\section{Parseval’s Theorem for DFT}

The energy in time and frequency domains is equal:

\[
\boxed{
\sum_{n=0}^{N-1}|x[n]|^2
=
\frac{1}{N}
\sum_{k=0}^{N-1}|X[k]|^2
}
\]

%%%%%%%%%%%%%%%%%%%%%%%%%%%%%%%%%%%%%%%%%%%%%%%%%%%%%%%%%%%%%%%%%%%%%%%%%%%%%%%%

\section{Example}

Let

\[
x[n]=\{1,1,1,1\}.
\]

Its DFT is

\[
X[k]=\{4,0,0,0\}.
\]

%%%%%%%%%%%%%%%%%%%%%%%%%%%%%%%%%%%%%%%%%%%%%%%%%%%%%%%%%%%%%%%%%%%%%%%%%%%%%%%%

\subsection{Time-Domain Energy}

\[
E_t=1+1+1+1=4.
\]

%%%%%%%%%%%%%%%%%%%%%%%%%%%%%%%%%%%%%%%%%%%%%%%%%%%%%%%%%%%%%%%%%%%%%%%%%%%%%%%%

\subsection{Frequency-Domain Energy}

\[
E_f=
\frac{1}{4}(16)=4.
\]

Hence Parseval’s theorem is verified.

%%%%%%%%%%%%%%%%%%%%%%%%%%%%%%%%%%%%%%%%%%%%%%%%%%%%%%%%%%%%%%%%%%%%%%%%%%%%%%%%

\section{Spectral Leakage}

The DFT assumes the observed sequence is one period of a periodic signal.

If the signal does not contain an integer number of cycles within the observation
window, energy spreads into adjacent bins.

This phenomenon is called \textbf{spectral leakage}.

%%%%%%%%%%%%%%%%%%%%%%%%%%%%%%%%%%%%%%%%%%%%%%%%%%%%%%%%%%%%%%%%%%%%%%%%%%%%%%%%

\section{Example of Leakage}

Suppose

\[
x[n]=\cos\left(\frac{2\pi(1.5)}{8}n\right),
\qquad N=8.
\]

The frequency lies between bins 1 and 2, so the DFT energy spreads across several
bins instead of appearing in a single bin.

%%%%%%%%%%%%%%%%%%%%%%%%%%%%%%%%%%%%%%%%%%%%%%%%%%%%%%%%%%%%%%%%%%%%%%%%%%%%%%%%

\section{Windowing}

To reduce leakage, the sequence is multiplied by a window.

%%%%%%%%%%%%%%%%%%%%%%%%%%%%%%%%%%%%%%%%%%%%%%%%%%%%%%%%%%%%%%%%%%%%%%%%%%%%%%%%

\subsection{Rectangular Window}

\[
w_R[n]=1.
\]

Narrow main lobe, high side lobes.

%%%%%%%%%%%%%%%%%%%%%%%%%%%%%%%%%%%%%%%%%%%%%%%%%%%%%%%%%%%%%%%%%%%%%%%%%%%%%%%%

\subsection{Hanning Window}

\[
\boxed{
w_H[n]=
0.5-0.5\cos\left(\frac{2\pi n}{N-1}\right)
}
\]

Wider main lobe, lower side lobes.

%%%%%%%%%%%%%%%%%%%%%%%%%%%%%%%%%%%%%%%%%%%%%%%%%%%%%%%%%%%%%%%%%%%%%%%%%%%%%%%%

\subsection{Hamming Window}

\[
\boxed{
w_{Ham}[n]=
0.54-0.46\cos\left(\frac{2\pi n}{N-1}\right)
}
\]

Very commonly used in practice.

%%%%%%%%%%%%%%%%%%%%%%%%%%%%%%%%%%%%%%%%%%%%%%%%%%%%%%%%%%%%%%%%%%%%%%%%%%%%%%%%

\section{Window Comparison}

\begin{table}[H]
\centering
\caption{Window characteristics}
\begin{tabular}{|p{4cm}|p{3cm}|p{4cm}|}
\hline
\textbf{Window} & \textbf{Main Lobe} & \textbf{Side Lobes}\\
\hline
Rectangular & Narrowest & Highest\\
\hline
Hanning & Wider & Lower\\
\hline
Hamming & Slightly wider & Much lower\\
\hline
\end{tabular}
\end{table}

%%%%%%%%%%%%%%%%%%%%%%%%%%%%%%%%%%%%%%%%%%%%%%%%%%%%%%%%%%%%%%%%%%%%%%%%%%%%%%%%

\section{Zero Padding}

Appending zeros to a sequence increases the DFT length.

%%%%%%%%%%%%%%%%%%%%%%%%%%%%%%%%%%%%%%%%%%%%%%%%%%%%%%%%%%%%%%%%%%%%%%%%%%%%%%%%

\subsection{Important Point}

Zero padding:

\begin{itemize}
\item does not increase actual frequency resolution,
\item only provides more interpolated spectral samples,
\item makes plots appear smoother.
\end{itemize}

%%%%%%%%%%%%%%%%%%%%%%%%%%%%%%%%%%%%%%%%%%%%%%%%%%%%%%%%%%%%%%%%%%%%%%%%%%%%%%%%

\section{Frequency Resolution}

For a sampling frequency \(f_s\),

\[
\boxed{
\Delta f=\frac{f_s}{N}
}
\]

To improve resolution:

\begin{itemize}
\item increase \(N\),
\item or observe the signal for a longer duration.
\end{itemize}

%%%%%%%%%%%%%%%%%%%%%%%%%%%%%%%%%%%%%%%%%%%%%%%%%%%%%%%%%%%%%%%%%%%%%%%%%%%%%%%%

\section{Solved Example 9.1}

Find the 4-point DFT of

\[
x[n]=\{1,1,1,1\}.
\]

%%%%%%%%%%%%%%%%%%%%%%%%%%%%%%%%%%%%%%%%%%%%%%%%%%%%%%%%%%%%%%%%%%%%%%%%%%%%%%%%

\subsection{Solution}

\[
X[0]=4.
\]

Using orthogonality,

\[
X[1]=X[2]=X[3]=0.
\]

\[
\boxed{
X[k]=\{4,0,0,0\}
}
\]

%%%%%%%%%%%%%%%%%%%%%%%%%%%%%%%%%%%%%%%%%%%%%%%%%%%%%%%%%%%%%%%%%%%%%%%%%%%%%%%%

\section{Solved Example 9.2}

Find the 4-point DFT of

\[
x[n]=\{1,0,-1,0\}.
\]

For \(N=4\),

\[
W_4=e^{-j\pi/2}=-j.
\]

%%%%%%%%%%%%%%%%%%%%%%%%%%%%%%%%%%%%%%%%%%%%%%%%%%%%%%%%%%%%%%%%%%%%%%%%%%%%%%%%

\subsection{Compute \(X[0]\)}

\[
X[0]=1+0-1+0=0.
\]

%%%%%%%%%%%%%%%%%%%%%%%%%%%%%%%%%%%%%%%%%%%%%%%%%%%%%%%%%%%%%%%%%%%%%%%%%%%%%%%%

\subsection{Compute \(X[1]\)}

\[
X[1]=1+0+(-1)(-1)+0=2.
\]

%%%%%%%%%%%%%%%%%%%%%%%%%%%%%%%%%%%%%%%%%%%%%%%%%%%%%%%%%%%%%%%%%%%%%%%%%%%%%%%%

\subsection{Compute \(X[2]\)}

\[
X[2]=1+0-1+0=0.
\]

%%%%%%%%%%%%%%%%%%%%%%%%%%%%%%%%%%%%%%%%%%%%%%%%%%%%%%%%%%%%%%%%%%%%%%%%%%%%%%%%

\subsection{Compute \(X[3]\)}

\[
X[3]=2.
\]

%%%%%%%%%%%%%%%%%%%%%%%%%%%%%%%%%%%%%%%%%%%%%%%%%%%%%%%%%%%%%%%%%%%%%%%%%%%%%%%%

\subsection{Result}

\[
\boxed{
X[k]=\{0,2,0,2\}
}
\]

%%%%%%%%%%%%%%%%%%%%%%%%%%%%%%%%%%%%%%%%%%%%%%%%%%%%%%%%%%%%%%%%%%%%%%%%%%%%%%%%

\section{Solved Example 9.3}

Find the circular convolution of

\[
x[n]=\{1,2,3\},
\qquad
h[n]=\{1,1,1\}.
\]

Using the circular convolution formula,

\[
y[0]=1+2+3=6,
\]

\[
y[1]=6,
\qquad
y[2]=6.
\]

Hence,

\[
\boxed{
y[n]=\{6,6,6\}
}
\]

%%%%%%%%%%%%%%%%%%%%%%%%%%%%%%%%%%%%%%%%%%%%%%%%%%%%%%%%%%%%%%%%%%%%%%%%%%%%%%%%

\section{Solved Example 9.4}

A sequence of length \(N=16\) is sampled at

\[
f_s=8\text{ kHz}.
\]

Find the frequency resolution.

%%%%%%%%%%%%%%%%%%%%%%%%%%%%%%%%%%%%%%%%%%%%%%%%%%%%%%%%%%%%%%%%%%%%%%%%%%%%%%%%

\subsection{Solution}

\[
\Delta f=
\frac{8000}{16}=500\text{ Hz}.
\]

\[
\boxed{
\Delta f=500\text{ Hz}
}
\]

%%%%%%%%%%%%%%%%%%%%%%%%%%%%%%%%%%%%%%%%%%%%%%%%%%%%%%%%%%%%%%%%%%%%%%%%%%%%%%%%

\section{Solved Example 9.5}

A 1 kHz sinusoid is sampled at 8 kHz and an 8-point DFT is computed.

%%%%%%%%%%%%%%%%%%%%%%%%%%%%%%%%%%%%%%%%%%%%%%%%%%%%%%%%%%%%%%%%%%%%%%%%%%%%%%%%

\subsection{Bin Frequency}

\[
f_k=\frac{k}{8}\times8000.
\]

Thus,

\[
f_1=1000\text{ Hz}.
\]

The sinusoid appears in

\[
\boxed{
k=1
}
\]

(and \(k=7\) due to conjugate symmetry).

%%%%%%%%%%%%%%%%%%%%%%%%%%%%%%%%%%%%%%%%%%%%%%%%%%%%%%%%%%%%%%%%%%%%%%%%%%%%%%%%

\section{Important GATE Points}

\begin{enumerate}
\item DFT assumes periodic extension.
\item Circular convolution results from periodicity.
\item Zero padding is required for linear convolution.
\item Real sequences exhibit conjugate symmetry.
\item Frequency resolution is \(f_s/N\).
\item Spectral leakage occurs when the signal frequency is not exactly at a DFT bin.
\item Windowing trades main-lobe width for lower side lobes.
\end{enumerate}

%%%%%%%%%%%%%%%%%%%%%%%%%%%%%%%%%%%%%%%%%%%%%%%%%%%%%%%%%%%%%%%%%%%%%%%%%%%%%%%%

\section{Chapter Summary}

In this chapter, we studied the Discrete Fourier Transform, its inverse transform,
twiddle factors and orthogonality.

We developed the concepts of circular convolution, linear convolution using DFT,
Parseval’s theorem, spectral leakage, windowing and zero padding.

These concepts form the mathematical foundation of FFT-based spectrum analyzers,
digital filters, OFDM systems and modern DSP processors.

%%%%%%%%%%%%%%%%%%%%%%%%%%%%%%%%%%%%%%%%%%%%%%%%%%%%%%%%%%%%%%%%%%%%%%%%%%%%%%%%

\begin{tcolorbox}[title=One-Minute Revision,colback=blue!5]

\[
X[k]=
\sum_{n=0}^{N-1}x[n]W_N^{kn}
\]

\[
x[n]=
\frac{1}{N}
\sum_{k=0}^{N-1}X[k]W_N^{-kn}
\]

\[
W_N=e^{-j2\pi/N}
\]

\[
y[n]=x[n]\circledast h[n]
\]

\[
x\circledast h
\xleftrightarrow{\text{DFT}}
XH
\]

\[
\Delta f=\frac{f_s}{N}
\]

\[
\sum |x[n]|^2=
\frac{1}{N}\sum |X[k]|^2
\]

Linear convolution:

\[
N\ge L_x+L_h-1
\]

Real sequence:

\[
X[N-k]=X^*[k]
\]

\end{tcolorbox}
%%%%%%%%%%%%%%%%%%%%%%%%%%%%%%%%%%%%%%%%%%%%%%%%%%%%%%%%%%%%%%%%%%%%%%%%%%%%%%%%